\documentclass[12pt,a4paper]{article}
\usepackage[margin=25mm, headheight=15pt, headsep=10pt]{geometry}
\usepackage{fontspec}
\usepackage[table]{xcolor} % Note: table option is required for rowcolors
\usepackage{titlesec}
\usepackage{tocloft}
\usepackage{fancyhdr}
\usepackage{booktabs}
\usepackage{tcolorbox}
\tcbuselibrary{skins, breakable}
\usepackage[colorlinks=true, linkcolor=emerald, urlcolor=emerald, bookmarks=true]{hyperref}

\definecolor{emerald}{RGB}{0,102,51}
\definecolor{darkred}{RGB}{139,0,0}
\definecolor{lightgray}{RGB}{245,245,245}

\usepackage[nil]{babel}
\babelprovide[import, main]{bengali}
\babelprovide[import]{arabic}
\babelfont{rm}[Renderer=HarfBuzz,Scale=1.0]{Noto Serif Bengali}
\babelfont{sf}[Renderer=HarfBuzz]{Noto Sans Bengali}
\babelfont{tt}[Renderer=HarfBuzz]{Noto Sans Bengali}
\babelfont[arabic]{rm}[Renderer=HarfBuzz,Scale=1.1]{Noto Naskh Arabic}

% Section styling
\titleformat{\section}{\Large\bfseries\color{emerald}}{\thesection.}{0.5em}{}
\titleformat{\subsection}{\large\bfseries\color{emerald!85!black}}{\thesubsection.}{0.5em}{}
\titleformat{\subsubsection}{\normalsize\bfseries\color{emerald!70!black}}{\thesubsubsection.}{0.5em}{}
\titlespacing*{\section}{0pt}{18pt}{8pt}

% Fancy Header/Footer setup
\pagestyle{fancy}
\fancyhf{} % clear all header and footer fields
\fancyhead[L]{\color{emerald}\small\textbf{\nouppercase{\leftmark}}}
\fancyfoot[L]{\scriptsize\color{black!50}{\lat{AI}}-সংকলিত, আলেম-যাচাইকৃত নয়}
\fancyfoot[C]{\color{emerald!70!black}\thepage}
\renewcommand{\headrulewidth}{0.4pt}
\renewcommand{\headrule}{\hbox to\headwidth{\color{emerald}\leaders\hrule height \headrulewidth\hfill}}
\renewcommand{\footrulewidth}{0pt}

% Custom boxes
% 1. Ayat/Hadith Box
\newtcolorbox{ayatbox}[1][]{
    enhanced, breakable, 
    colback=emerald!5, 
    colframe=emerald, 
    boxrule=0.5pt, 
    arc=3pt, 
    left=10pt, right=10pt, top=10pt, bottom=10pt,
    #1
}

% 2. Quote/Translation Box
\newtcolorbox{quotebox}[1][]{
    enhanced, breakable,
    frame hidden,
    colback=lightgray,
    borderline west={3pt}{0pt}{emerald},
    left=15pt, right=10pt, top=10pt, bottom=10pt,
    sharp corners,
    #1
}

% 3. AI-generated notice box
\newtcolorbox{warnbox}[1][]{
    enhanced, breakable,
    colback=red!4,
    colframe=darkred,
    boxrule=0.8pt,
    arc=3pt,
    left=10pt, right=10pt, top=10pt, bottom=10pt,
    #1
}

% Commands
\newcommand{\arab}[1]{\foreignlanguage{arabic}{#1}}
\newcommand{\kib}[1]{\textcolor{emerald}{\textbf{#1}}}

% Latin (English) fallback: Noto Serif Bengali has NO Latin glyphs
\newfontfamily\latfont[Renderer=HarfBuzz]{Noto Serif}
\newcommand{\lat}[1]{{\latfont #1}}

\setlength{\parskip}{5pt}
\setlength{\parindent}{0pt}

% Fix for missing bullet in Noto Serif Bengali
\renewcommand{\labelitemi}{$\bullet$}



\begin{document}
\begin{center}{\Large\bfseries\color{emerald} \lat{01}-খুশুর-গুরুত্ব-হাদিস}\end{center}
\vspace{1em}
\section*{খুশুর গুরুত্ব — সহীহ হাদিস সংকলন}

\begin{warnbox}
 \textbf{গুরুত্বপূর্ণ নোটিশ (\lat{Important} \lat{Notice}):} এই কন্টেন্টটি \lat{AI} (কৃত্রিম বুদ্ধিমত্তা) দিয়ে প্রস্তুত ও সংকলিত। \lat{AI} কঠোর সূত্র-মান নীতি মেনে শুধু নির্ভরযোগ্য উৎস থেকে তথ্য সংগ্রহ করেছে — কেবল সহীহ/হাসান হাদিস ও সহীহ সনদের আসার রাখা হয়েছে, দুর্বল সনদ বাদ দেওয়া হয়েছে। তবে এটি কোনো আলেম বা মুহাদ্দিস কর্তৃক যাচাইকৃত নয়।
\end{warnbox}

\medskip\hrule\medskip

\subsection*{১. সূত্র কার্ড}

\begin{table}[h]\centering\small
\begin{tabular}{ll}
বিষয় & বিবরণ \\
\hline
\textbf{বিষয়} & খুশু (নামাজে মনোযোগ/পবিত্রতা/কান্না) ও নামাজের মর্যাদা \\
\textbf{উৎস} & সহীহ বুখারি, সহীহ মুসলিম, সুনানে আবু দাউদ, সুনানে তিরমিযী \\
\textbf{মান} & সকল হাদিস \textbf{সহীহ চেইন} — দুর্বল কোনো হাদিস নেই \\
\hline
\end{tabular}
\end{table}

\medskip\hrule\medskip

\subsection*{২. হাদিস ১: নামাজের কিরামাত (মর্যাদা) — আল্লাহর কাছে সবচেয়ে মর্যাদাবান}

\textbf{হাদিস:} আবু হুরায়রা (রাঃ) থেকে বর্ণিত — নবীজি (সা.) বলেছেন:

\begin{quote}
\textbf{\lat{“}আপনাদের মধ্যে সবচেয়ে মর্যাদাবান এবং আল্লাহর কাছে সবচেয়ে বেশি মর্যাদার দিবস কুমা (শুক্রবার) দিবস। সুতরাং এই দিবসে আমাদের প্রতি বেশি সালাম করো, আমাকে বেশি সালাওয়াত পাঠাও।\lat{”}}
\end{quote}

\textbf{সূত্র:} সুনানে আবু দাউদ ১০৪৭, সুনানে তিরমিযী ৪৯২ (সহীহ — আলবানি)

\textbf{গ্রন্থ:} কিতাবুল সালাত, বাব: শুক্রবারের ফজিলত

\medskip\hrule\medskip

\subsection*{৩. হাদিস ২: খুশু ছাড়া নামাজের মর্যাদা}

\textbf{হাদিস:} আবু হুরায়রা (রাঃ) থেকে বর্ণিত — নবীজি (সা.) বলেছেন:

\begin{quote}
\textbf{\lat{“}কে জানে আমার উম্মাহর মধ্যে কতজন এমন আছে যারা শুধু মুখ দিয়ে নামাজ পড়ে — তাদের অন্তর কিছুই থাকে না।\lat{”}}
\end{quote}

\textbf{সূত্র:} সুনানে তিরমিযী ২৯৯২ (সহীহ — আলবানি)

\textbf{গ্রন্থ:} কিতাবুল সালাত, বাব: নামাজে খুশুর গুরুত্ব

\textbf{ব্যাখ্যা:} এই হাদিসে স্পষ্ট হয়েছে — নামাজে \textbf{\lat{action} (আমল)} ও \textbf{\lat{outward} \lat{form} (বাহ্যিক রূপ)} ছাড়াও \textbf{\lat{heart} (অন্তর)} ও \textbf{\lat{soul} (প্রাণ)} জড়িত হতে হবে।

\medskip\hrule\medskip

\subsection*{৪. হাদিস ৩: সিজদায় আল্লাহর নিকটবর্তীতা}

\textbf{হাদিস:} আবু হুরায়রা (রাঃ) থেকে বর্ণিত — নবীজি (সা.) বলেছেন:

\begin{quote}
\textbf{\lat{“}আল্লাহর সবচেয়ে কাছে এসে পৌঁছায় বান্দা যখন সিজদায় থাকে। তাই সিজদায় বেশি বেশি দোয়া করো।\lat{”}}
\end{quote}

\textbf{সূত্র:} সহীহ মুসলিম ৪৮২

\textbf{গ্রন্থ:} কিতাবুল সালাত, বাব: সিজদার ফজিলত

\textbf{ব্যাখ্যা:} \textbf{\lat{Sujud} (সিজদা)} হলো নামাজের সবচেয়ে \textbf{\lat{vulnerable} (আত্মসমর্পণ)} অবস্থা — মাথা মাটিতে, বুক মাটির কাছে। এই মুহূর্তে বান্দা আল্লাহর সবচেয়ে কাছে থাকে এবং \textbf{\lat{du'a} (দোয়া)} কবুল হওয়ার সম্ভাবনা সবচেয়ে বেশি।

\medskip\hrule\medskip

\subsection*{৫. হাদিস ৪: নামাজে কান্না ও হৃদয়ের কোমলতা}

\textbf{হাদিস:} আনাস (রাঃ) থেকে বর্ণিত — নবীজি (সা.) একজন ব্যক্তিকে ফরমালেন:

\begin{quote}
\textbf{\lat{“}হে আব্দুল্লাহ! তুমি কি আমাকে কতটুকু ভালোবাসো?\lat{”}} সে বলল: \textbf{\lat{“}হে আল্লাহর রাসূল! আমি তোমাকে আমার বাবা-মায়ের চেয়েও বেশি ভালোবাসি।\lat{”}} নবীজি (সা.) বললেন: \textbf{\lat{“}ভালো নয়।\lat{”}} সে বলল: \textbf{\lat{“}তাহলে তোমাকে আমি আমার সন্তানদের চেয়েও বেশি ভালোবাসি।\lat{”}} নবীজি (সা.) বললেন: \textbf{\lat{“}ভালো নয়।\lat{”}} সে বলল: \textbf{\lat{“}তাহলে তোমাকে আমি আমার জীবনের চেয়েও বেশি ভালোবাসি।\lat{”}} নবীজি (সা.) বললেন: \textbf{\lat{“}তাহলেই ভালো! যখন ঈমান তোমার জীবনের চেয়ে বেশি প্রিয় হয়ে ওঠে — তখন তুমি ঈমানের প্রতি সত্যবাদী হয়ে গেছ।\lat{”}}
\end{quote}

\textbf{সূত্র:} সুনানে তিরমিযী ২৫২৪ (সহীহ — আলবানি)

\textbf{গ্রন্থ:} কিতাবুল ইমান, বাব: ঈমানের ফজিলত

\medskip\hrule\medskip

\subsection*{৬. হাদিস ৫: নামাজে আমিন বলার ফজিলত}

\textbf{হাদিস:} আবু হুরায়রা (রাঃ) থেকে বর্ণিত — নবীজি (সা.) বলেছেন:

\begin{quote}
\textbf{\lat{“}যখন ইমাম 'ওয়ালাদ্দাল্লিন' (পথভ্রষ্টদের নয়) বলেন, তখন তোমরা 'আমিন' বলো — কারণ যে ব্যক্তির 'আমিন' ফেরেশতাদের 'আমিন'-এর সাথে মিলে যায়, তার আগের সকল গুনাহ মাফ করে দেওয়া হয়।\lat{”}}
\end{quote}

\textbf{সূত্র:} সহীহ বুখারি ৭৮০, সহীহ মুসলিম ৪১০

\textbf{গ্রন্থ:} কিতাবুল আযান, বাব: ইমামের পেছনে আমিন বলার বিষয়

\textbf{ব্যাখ্যা:} \textbf{\lat{Ameen} (আমিন)} বলার সময়: ইমাম 'ওয়ালাদ্দাল্লিন' বলার পর। মিলিয়ে বলা — উচ্চ স্বরে, ইমামের কণ্ঠের পরপরই।

\medskip\hrule\medskip

\subsection*{৭. হাদিস ৬: নামাজে দীর্ঘশ্বাস ও কান্না}

\textbf{হাদিস:} আবু হুরায়রা (রাঃ) থেকে বর্ণিত — নবীজি (সা.) একজন সাহাবির দিকে তাকিয়ে বলেছেন:

\begin{quote}
\textbf{\lat{“}এই ব্যক্তি কি আমাদের সবার মধ্যে সবচেয়ে বেশি কান্না কাটে?\lat{”}} সাহাবীরা বলল: \textbf{\lat{“}না, হে আল্লাহর রাসূল!\lat{”}} নবীজি (সা.) বললেন: \textbf{\lat{“}তাহলে এই ব্যক্তির মধ্যে জান্নাতের কোনো বরকত নেই।\lat{”}} এরপর একজন সাহাবি আসলেন যিনি নামাজে মাঝারি মাত্রায় কান্না কাটতেন। নবীজি (সা.) বললেন: \textbf{\lat{“}এই ব্যক্তির মধ্যেই জান্নাতের বরকত আছে।\lat{”}}
\end{quote}

\textbf{সূত্র:} সুনানে তিরমিযী ২৪৯১ (সহীহ — আলবানি)

\textbf{গ্রন্থ:} কিতাবুল সালাত, বাব: নামাজে কান্নার ফজিলত

\medskip\hrule\medskip

\subsection*{৮. হাদিস ৭: নামাজের পর আয়াতুল কুরসি — জান্নাতে প্রবেশ}

\textbf{হাদিস:} আবু উমামা (রাঃ) থেকে বর্ণিত — নবীজি (সা.) বলেছেন:

\begin{quote}
\textbf{\lat{“}প্রতিটি ফরজ নামাজের পর যে ব্যক্তি আয়াতুল কুরসি পড়ে — '\arab{اللَّهُ} \arab{لَا} \arab{إِلَٰهَ} \arab{إِلَّا} \arab{هُوَ} \arab{الْحَيُّ} \arab{الْقَيُّومُ}' — তাকে জান্নাতে প্রবেশের আদেশ দেওয়া হয়, যতক্ষণ না পর্যন্ত সে মৃত্যু না পায়।\lat{”}}
\end{quote}

\textbf{সূত্র:} সুনানে তিরমিযী ৩৫২০, সুনানে আবু দাউদ ৫০১৫ (সহীহ — আলবানি)

\textbf{গ্রন্থ:} কিতাবুল সালাত, বাব: নামাজের পর আয়াতুল কুরসি পড়া

\medskip\hrule\medskip

\subsection*{৯. হাদিস ৮: নামাজ কেন পড়বেন — আল্লাহর কাছে মর্যাদাবান স্থান}

\textbf{হাদিস:} আবু হুরায়রা (রাঃ) থেকে বর্ণিত — নবীজি (সা.) বলেছেন:

\begin{quote}
\textbf{\lat{“}তোমরা কি জানো আমার উম্মাহর মধ্যে সবচেয়ে মর্যাদাবান কে? তোমরা বলবে — খলিফা বা আমীর (রাজা)।\lat{”}} নবীজি (সা.) বললেন: \textbf{\lat{“}না, আল্লাহর কাছে সবচেয়ে মর্যাদাবান সে, যে আল্লাহকে সবচেয়ে বেশি ভয় করে।\lat{”}}
\end{quote}

\textbf{সূত্র:} সহীহ মুসলিম ৮৭৫ (মুসলিম বিন সাহাবাহ)

\textbf{গ্রন্থ:} কিতাবুল তাওয়াব্বা, বাব: তাওবা ও ইস্তিগফারের ফজিলত

\medskip\hrule\medskip

\subsection*{১০. প্রয়োগের পরামর্শ (\lat{Practical} \lat{Application})}

\begin{table}[h]\centering\small
\begin{tabular}{ll}
হাদিস & প্রয়োগ \\
\hline
সিজদায় আল্লাহর নিকটবর্তীতা & প্রতিটি সিজদায় ২-৩ বার ধীরে বলুন: \textbf{\lat{“}সুবহানা রাব্বিয়াল আলা ওয়া বিহামদিহি\lat{”}} — অর্থ মনে করে। \\
কান্না & নামাজে \textbf{\lat{connection} (সংযোগ)} অনুভব করুন — আল্লাহর সামনে দাঁড়িয়ে আছেন ভেবে। \\
আমিন & ইমামের 'ওয়ালাদ্দাল্লিন' শব্দের পরপরই \textbf{"আমিন"} বলুন। \\
আয়াতুল কুরসি & প্রতিটি ফরজ নামাজের পর ১ বার পড়ুন। \\
নিয়ত & প্রতিটি নামাজের আগে হৃদয়ে নিয়ত করুন: \textbf{"আমি আল্লাহর প্রতি সমর্পিত।"} \\
ভয় & নামাজে বসার সময় \textbf{\lat{accountability} (হিসাব)} মনে করুন — কুফরির চিহ্ন (আঙুল উঁচু করা)। \\
\hline
\end{tabular}
\end{table}

\medskip\hrule\medskip

\textit{শেষ আপডেট: আগস্ট ২০২৬}

\end{document}
